\chapter{Conclusion and Future Work}
\label{ch:conclusion}

\section{Conclusion}

This thesis presented a quantization and hardware validation workflow connecting Python export, C++ reference inference, Python hardware-like simulation, and RTL testbench validation. The C++ reference sweep showed that int16 inference largely preserves float model behavior, while some layers, particularly patch embedding and input projection, are sensitive to int8 quantization.

Stage-level comparison between \texttt{cpp\_like} and \texttt{hw\_like} traces identified the selective scan as the dominant source of hardware-oriented fixed-point error in case 04. A scaled-state representation reduced the average selective scan MAE from 0.121098 to 0.090323, corresponding to a 25.41\% improvement.

\section{Future Work}

Future work includes:

\begin{itemize}
  \item Extending the scaled-state method to more datasets and model configurations.
  \item Automating scale selection across more stages.
  \item Improving RTL coverage for the full front-end, including patch embedding.
  \item Evaluating resource, timing, and energy trade-offs of mixed int8/int16 configurations.
  \item Integrating hardware-in-the-loop validation on FPGA.
\end{itemize}

